\section{结果读取说明}

本章计算目的：说明如何读取表 5 和表 6。表 5 是目标相关源项的 100--1000 m 客观参考；表 6 是当前变频器--水泵系统实验背景的 0--5 m 等效场强参考。两张表不合并排序。

\subsection{目标源项幅值范围读取}

在表 5 的参数包络内，UEP/ICCP 静态或慢变电场在 100--1000 m 的 reference 量级为 $5.97\times10^{-6}$ 到 $5.97\times10^{-9}\ \mathrm{V/m}$，包络范围为 $3.98\times10^{-7}$ 到 $1.59\times10^{-4}\ \mathrm{V/m}$（100 m）和 $3.98\times10^{-10}$ 到 $1.59\times10^{-7}\ \mathrm{V/m}$（1000 m）。轴频基频 reference 量级为 $1.79\times10^{-7}$ 到 $1.79\times10^{-10}\ \mathrm{V/m}$，二/三阶谐波 reference 量级低于基频。

尾流/流动诱导电场在尾流中心线模型下的 reference 量级为 $1.98\times10^{-6}$ 到 $1.25\times10^{-7}\ \mathrm{V/m}$。该数值对应下游中心线距离 $x$，不是径向远场偶极距离 $R$。

\subsection{频率特征读取}

UEP/ICCP 静态或慢变电场对应 0 Hz 或低频慢变分量。
轴频基频具有 $f_s$ 线谱，典型参考频率单独按 3 Hz 读取；二/三阶谐波具有 $2f_s/3f_s$ 线谱，典型参考为 6/9 Hz。
尾流/流动诱导电场按低频或宽带慢变读取。变频器--水泵系统高频电磁背景按 VFD 输出基波 $f_{\mathrm{out}}$、基波谐波、PWM 载波 $f_c$ 及谐波、旁带 $nf_c\pm mf_{\mathrm{out}}$、50--200 kHz 高频振铃/共振以及 50 Hz 输入侧附加背景读取。

\subsection{模型可信度读取}

UEP/ICCP 与轴频行使用相同的远场偶极框架，模型边界清晰，但数值受 $P_0$、$\sigma$、$C_\theta$ 和调制深度影响。尾流模型使用中心线速度衰减和地磁耦合估算，模型可信度低于 UEP/轴频远场模型。变频器--水泵系统高频电磁背景不作为真实水下目标源项远场预测，而作为当前实验系统背景场。

\subsection{变频器--水泵系统高频电磁背景读取}

变频器--水泵系统高频电磁背景不作为真实水下目标源项远场预测，而作为当前实验系统背景场。按频率相关共模/漏电流等效偶极模型估算，若某一频率处存在 $100\ \mu\mathrm{A}$ 级可见共模/漏电流，则 1--3 m 实验距离处背景约为 $\mu\mathrm{V/m}$ 到十几 $\mu\mathrm{V/m}$ 量级；若达到 1 mA 级较强耦合，则可达到数十至数百 $\mu\mathrm{V/m}$ 量级。该表用于和表 5 中目标源项强度进行实验背景对照。

50--100 kHz 和 100--200 kHz 显著高于 UEP/ICCP、轴频、尾流的主要频率范围，不宜解释为真实水下目标主源项。该频段更可能来自 VFD PWM 开关噪声、快速电压边沿、电机电缆共模电流、泵壳/水体/接地耦合、DAQ/MEMS 高频响应。

\subsection{中立性约束的读取方式}

表 5 和表 6 只用于客观参考：可以读取哪类目标源项幅值范围更高、哪类源项频率特征更清晰、哪类模型可信度较低，以及当前 VFD--水泵系统背景可能处于什么量级。表中不包含实验排序，也不根据强度或模型可信度替用户决定实验取舍。

本章对客观参考表的作用：说明表 5 和表 6 的分工，使目标源项远场估算与 VFD--水泵系统实验背景对照保持分离。
